\SourceLeaf{082}%
servationes complete compensatae omnibus aequationibus conditionalibus \(X=0,\allowbreak Y=0,\allowbreak Z=0\) etc. satisfaciunt, incomplete compensatae vero vel nullis vel saltem non omnibus: compensatio itaque, per quam omnibus aequationibus conditionalibus satisfit, necessario est ipsa completa.

\Article{19}
Iam quum ex ipsa notione compensationis sponte sequatur, aggregata duarum compensationum iterum constituere compensationem, facile perspicitur, nihil interesse, utrum praecepta, per quae compensatio perfecta eruenda est, immediate ad observationes primitivas applicentur, an ad observationes incomplete iam compensatas.

Revera constituant \(-\theta,-\theta',-\theta''\) etc. systema compensationis incompletae, quod prodierit e formulis (I)
\[
\begin{aligned}
\theta p&=A^0a+B^0b+C^0c+\text{etc.}\\
\theta'p'&=A^0a'+B^0b'+C^0c'+\text{etc.}\\
\theta''p''&=A^0a''+B^0b''+C^0c''+\text{etc.}\\
&\text{etc.}
\end{aligned}
\]

Quum observationes his compensationibus mutatae omnibus aequationibus conditionalibus non satisfacere supponantur, sint \(\mathfrak A^*,\mathfrak B^*,\mathfrak C^*\) etc. valores, quos \(X,Y,Z\) etc. ex illarum substitutione nanciscuntur. Quaerendi sunt numeri \(A^*,B^*,C^*\) etc. aequationibus (II) satisfacientes
\[
\begin{aligned}
\mathfrak A^*&=A^*[aa]+B^*[ab]+C^*[ac]+\text{etc.}\\
\mathfrak B^*&=A^*[ab]+B^*[bb]+C^*[bc]+\text{etc.}\\
\mathfrak C^*&=A^*[ac]+B^*[bc]+C^*[cc]+\text{etc.}
\end{aligned}
\]
etc., quo facto compensatio completa\textperiodcentered\ observationum isto modo mutatarum efficitur per mutationes novas \(-\kappa,-\kappa',-\kappa''\) etc., ubi \(\kappa,\kappa',\kappa''\) etc. computandae sunt per formulas (III)
\[
\begin{aligned}
\kappa p&=A^*a+B^*b+C^*c+\text{etc.}\\
\kappa'p'&=A^*a'+B^*b'+C^*c'+\text{etc.}\\
\kappa''p''&=A^*a''+B^*b''+C^*c''+\text{etc.}
\end{aligned}
\]
etc. Iam inquiramus, quomodo hae correctiones cum compensatione completa
\SourceLeaf{083}
observationum primitivarum cohaereant. Primo manifestum est, haberi
\[
\begin{aligned}
\mathfrak A^*&=\mathfrak A-a\theta-a'\theta'-a''\theta''-\text{etc.}\\
\mathfrak B^*&=\mathfrak B-b\theta-b'\theta'-b''\theta''-\text{etc.}\\
\mathfrak C^*&=\mathfrak C-c\theta-c'\theta'-c''\theta''-\text{etc.}
\end{aligned}
\]
etc. Substituendo in his aequationibus pro \(\theta,\theta',\theta''\) etc. valores ex (I), nec non pro \(\mathfrak A^*,\mathfrak B^*,\mathfrak C^*\) etc. valores ex II, invenimus
\[
\begin{aligned}
\mathfrak A&=(A^0+A^*)[aa]+(B^0+B^*)[ab]+(C^0+C^*)[ac]+\text{etc.}\\
\mathfrak B&=(A^0+A^*)[ab]+(B^0+B^*)[bb]+(C^0+C^*)[bc]+\text{etc.}\\
\mathfrak C&=(A^0+A^*)[ac]+(B^0+B^*)[bc]+(C^0+C^*)[cc]+\text{etc.}
\end{aligned}
\]
etc., unde patet, correlata aequationum conditionalium aequationibus (12) satisfacientia esse
\[
A=A^0+A^*,\quad B=B^0+B^*,\quad C=C^0+C^*\ \text{etc.}
\]
Hinc vero aequationes (10), I et III docent, esse
\[
\varepsilon=\theta+\kappa,\quad\varepsilon'=\theta'+\kappa'\quad\varepsilon''=\theta''+\kappa''\ \text{etc.}
\]
i. e. compensatio observationum perfecta eadem prodit, sive immediate computetur, sive mediate proficiscendo a compensatione manca.

\Article{20}
Quoties multitudo aequationum conditionalium permagna est, determinatio correlatorum \(A,B,C\) etc. per eliminationem directam tam prolixa evadere potest, ut calculatoris patientia ei impar sit: tunc saepenumero commodum esse poterit, compensationem completam per approximationes successivas adiumento theorematis art. praec. eruere. Distribuantur aequationes conditionales in duas pluresve classes, investigeturque primo compensatio, per quam aequationibus primae classis satisfit, neglectis reliquis. Dein tractentur observationes per hanc compensationem mutatae ita, ut solarum aequationum secundae classis ratio habeatur. Generaliter loquendo applicatio secundi compensationum systematis consensum cum aequationibus primae classis turbabit; quare, si duae tantummodo classes factae sunt, ad aequationes primae classis revertemur, tertiumque systema, quod huic satisfaciat, eruemus; dein observationes ter correctas compensationi quartae
